\documentclass[10pt,a4paper]{amsart}
\usepackage[margin=19mm]{geometry}
\usepackage{amsmath,amssymb,mathtools}
\usepackage[scr=rsfs,cal=euler]{mathalfa}
\usepackage{xfrac}
\usepackage[unicode,hidelinks,bookmarks=false]{hyperref}
\newcommand{\ie}{i.e.\ }
\newcommand{\cf}{cf.\ }
\newcommand{\resp}{respectively\ }
\newcommand{\Subsection}{\subsection}
\providecommand{\nobreakdash}{\nobreak\hbox{-}}
\setlength{\emergencystretch}{2em}
\multlinegap0pt
\usepackage{amssymb}
\usepackage{tikz}
\usepackage[shortlabels]{enumitem}
\usepackage[normalem]{ulem}
\newtheorem{lemma}{Lemma}
\newcommand{\im}{\operatorname{Im}}
\newcommand{\re}{\operatorname{Re}}
\title{The chiral gyrating H'-T surface family:\\ construction from the dual \texttt{qtz}--\texttt{qzd} nets and\\ existence proof using a toroidal Weierstrass method}
\author{Hao Chen, Shashank G. Markande, Matthias Saba, Gerd E. Schr\"oder-Turk, Elisabetta A. Matsumoto}
\date{Source: arXiv:2512.18308v1 (CC BY 4.0)}
\begin{document}
\maketitle

\noindent\emph{Source and licence.} The chiral gyrating H'-T surface family:\\ construction from the dual \texttt{qtz}--\texttt{qzd} nets and\\ existence proof using a toroidal Weierstrass method. Version arXiv:2512.18308v1 (CC BY 4.0), distributed by arXiv under the Creative Commons Attribution 4.0 International licence, http://creativecommons.org/licenses/by/4.0/, as recorded in the arXiv licence metadata for this exact version and shown on the abstract page https://arxiv.org/abs/2512.18308v1. Full licence notice: \texttt{licenses/arxiv-2512.18308-CC-BY-4.0.txt}.\par\smallskip

\noindent\textit{Editorial scope.} Counted unit: the article's lemma lemma (source0.tex lines 881--887). The article's complete original proof of it is delivered with it (source0.tex lines 889--950, one \texttt{proof} environment of 62 lines). No mathematics is written, repaired or re-derived: the statement and the proof are the article's own lines, in the article's own notation, and no retained line is substituted or re-flowed.  No further source-local context is needed: every cross-reference of every retained line resolves inside the two delivered spans.  Nothing was omitted from any retained span, so the package carries no omission record.  The retained lines cite 1 works, whose entries are transcribed verbatim from the article's own bibliography source in the e-print and delivered with short editorial labels recorded in the item's reference mapping.  No retained line contains a figure environment, an image-inclusion command or drawing code, so no image is delivered and the package declares no asset dependency.  Editorial formatting only, in addition: an AMS \texttt{amsart} preamble that restates the article's own declarations and macros used by the retained lines, namely the environments \texttt{lemma} on separate counters, together with the standard packages those lines need.  The retained environments print in this document as Lemma 1 in order of appearance, whereas the article prints the same items with the numbering of its own full document.  The complete native source of the article as distributed in the arXiv e-print for this exact version is bundled unchanged as \texttt{sources/191376/source0.tex}.
\begin{lemma}
	\[
		\lim_{\im\tau \to +\infty} \arg \psi(\tau)
		= \lim_{\im\tau \to +\infty} \arg \int_0^{\tau/3} G(z) dz
		= \frac{1}{2} - \frac{\re\tau}{3}.
	\]
\end{lemma}

\begin{proof}
	In the limit $\im\tau \to +\infty$, by the expansion~\cite[(20.2.1)]{DLMF}
	\[
		\theta(z;\tau) = 2 q^{1/4} \sin(z) + o(q^{1/4}),
	\]
	where $q = \exp(i \pi \tau)$, we have
	\begin{align*}
		G(z)
		&\sim
		\sqrt{-\frac{\sin(\pi 3z)}{\sin(\pi (3z - \tau)} e^{i \pi (2z - \tau/3)}}\\
			&= \sqrt{-\frac{e^{i \pi 3z} - e^{-i \pi 3z}}{e^{i \pi (3z - \tau)} - e^{-i \pi (3z - \tau)}}} e^{i \pi (z - \tau/6)},
		\end{align*}
		as $\im\tau \to +\infty$.

		Now we divide $\psi(\tau)$ into two integrals
		\[
			\psi(\tau) = \psi_1(\tau) + \psi_2(\tau) = \int_0^{\tau/6} G(z) dz + \int_{\tau/6}^{\tau/3} G(z) dz
		\]

		In the integral $\psi_1(\tau)$, $\tau - 3z > \tau/2$.  So as $\im \tau \to +\infty$, we 
		\begin{align}
			\psi_1(\tau)
			&\sim \int_0^{\tau/6} \sqrt{\frac{e^{i \pi 3z} - e^{-i \pi 3z}}{e^{-i \pi (3z - \tau)}}} e^{i \pi (z - \tau/6)} dz\nonumber\\
			&= e^{i \pi \tau/3} \int_0^{\tau/6} \sqrt{e^{i \pi 6z} - 1} e^{i \pi z} dz\nonumber\\
			&= \frac{i}{\pi} e^{i \pi \tau/3} \int_1^{e^{-i\pi\tau/6}} \frac{\sqrt{w^6-1}}{w^2} dw  & (w = e^{-i\pi z}).\label{eq:app1}
		\end{align}
		Note that
		\[
			\Big| \frac{\sqrt{w^6-1}}{w^2} \Big| < |w|,
		\]
		we have
		\begin{align*}
			| \text{\eqref{eq:app1}} |
			&\le \frac{1}{\pi} e^{-\pi \im\tau/3} \int_1^{e^{\pi\im\tau/6}} x dx \\
			& = \frac{1}{2\pi} (1 - e^{-\pi\im\tau/3}) \to \frac{1}{2\pi}
		\end{align*}
		as $\im\tau \to +\infty$.  So $\psi_1(\tau)$ is bounded.

		In the integral $\psi_2(\tau)$, $3z > \tau/2$.  So as $\im \tau \to +\infty$, we 
		\begin{align*}
			\psi_2(\tau)
			&\sim \int_{\tau/6}^{\tau/3} \sqrt{\frac{e^{-i \pi 3z}}{e^{i \pi (3z - \tau)} - e^{-i \pi (3z - \tau)}}} e^{i \pi (z - \tau/6)} dz\\
			&= \int_{0}^{\tau/6} \sqrt{\frac{e^{-i \pi (\tau - 3\zeta)}}{e^{- i \pi 3\zeta} - e^{i \pi 3\zeta}}} e^{i \pi (\tau/6 - \zeta)} dz & (3\zeta = \tau - 3z)\\
			&= e^{-i \pi \tau/3} \int_0^{\tau/6} \frac{e^{-i\pi \zeta}}{\sqrt{e^{-i\pi 6\zeta} -1}} d\zeta\\
			&= \frac{i}{\pi} e^{-i \pi \tau/3} \int_1^{e^{-i\pi \tau/6}} \frac{du}{\sqrt{u^6 -1}} & (u = e^{-i \pi \zeta})\\
			&= \frac{i}{\pi} e^{-i \pi \tau/3} \int_1^{e^{i\pi \tau/6}} \frac{-v dv}{\sqrt{1 - v^6}} & (v = 1/w)
		\end{align*}
		Note that the integral in the last line converge to the bounded real integral
		\[
			I = \int_0^1 \frac{x dx}{\sqrt{1 - x^6}} \in \mathbb{R}
		\]
		as $\im\tau \to +\infty$, so
		\[
			\psi_2(\tau) \sim \frac{i}{\pi} I e^{-i\pi \tau/3}.
		\]

		As $\im \tau \to +\infty$, $\psi_2$ is dominant, so
		\[
			\lim_{\im\tau \to +\infty} \arg \psi(\tau) =
			\lim_{\im\tau \to +\infty} \arg \psi_2(\tau) = \frac{1}{2} - \frac{\tau}{3}.
		\]
	\end{proof}

\begin{thebibliography}{99}
\bibitem[DLMF]{DLMF} {\it NIST Digital Library of Mathematical Functions}. \newblock \texttt{http://dlmf.nist.gov/}, Release 1.0.19 of 2018-06-22. \newblock F.~W.~J. Olver, A.~B. {Olde Daalhuis}, D.~W. Lozier, B.~I. Schneider, R.~F. Boisvert, C.~W. Clark, B.~R. Miller and B.~V. Saunders, eds.
\end{thebibliography}
\end{document}
